\chapter{Riferimento operativo}
\label{chap:reference}

I comandi riportati in questa appendice devono essere eseguiti dalla root del
repository.

\section{Ambiente e verifica}

\begin{longtable}{P{0.48\textwidth}P{0.40\textwidth}}
\toprule
\textbf{Comando} & \textbf{Uso}\\
\midrule\endhead
\cli{conda env create -f environment.yml} & Crea l'ambiente principale dal file versionato.\\
\cli{conda activate graph_rec} & Attiva l'ambiente prima degli entrypoint.\\
\cli{conda env update -n graph_rec -f environment.yml --prune} & Riallinea l'ambiente da \filepath{environment.yml}.\\
\cli{python -c "from graph_reconstruction.config import IMAGES_DIR, WORK_DIR; print(IMAGES_DIR, WORK_DIR)"} & Verifica i percorsi caricati da \filepath{.env}.\\
\bottomrule
\end{longtable}

\section{Comandi principali}

\begin{longtable}{P{0.43\textwidth}P{0.45\textwidth}}
\toprule
\textbf{Comando} & \textbf{Risultato}\\
\midrule\endhead
\cli{python run.py 7} & Ricostruisce soltanto il nodo \code{7}.\\
\cli{python run.py 2 7 15_1} & Ricostruisce gli ID indicati.\\
\cli{python run.py} & Ricostruisce tutti gli ID scoperti.\\
\cli{python evaluate.py 7} & Stampa le metriche del nodo e aggiorna il confronto con gli standard.\\
\cli{python evaluate.py} & Valuta tutti i risultati presenti.\\
\cli{python visualize.py 7} & Genera il plot, il GeoJSON sintetico e il report visuale del nodo.\\
\cli{python visualize.py --all} & Visualizza tutti i grafi disponibili.\\
\cli{python json_adapter_to_qgis.py} & Esporta in batch i grafi disponibili. Eseguire per ultimo se serve una geometria non \code{fake}.\\
\cli{python -m graph_reconstruction.visualization.shurikode_symbols --all} & Genera i render SVG con i simboli Shurikode.\\
\bottomrule
\end{longtable}

\manualnote[Wrapper Windows]{\filepath{tools/run_pipeline.cmd} concatena run,
evaluate e visualize, ma preferisce \code{py -3} quando il launcher è presente.
Può quindi bypassare Conda e stampa percorsi finali basati sulla directory
\filepath{output/}. Con una \envvar{WORK_DIR} esterna è preferibile eseguire i
tre comandi \code{python} in modo esplicito.}

\section{Percorsi e artefatti}

\begin{longtable}{P{0.43\textwidth}P{0.45\textwidth}}
\toprule
\textbf{Percorso} & \textbf{Contenuto}\\
\midrule\endhead
\filepath{<IMAGES_DIR>/node<id>_<frame>.jpg} & Immagini della lavorazione; la ricerca non è ricorsiva.\\
\filepath{<LABELS_EXT_DIR>/<image_stem>_extended.txt} & Label estese usate dalla valutazione.\\
\filepath{<ANNOTATIONS_DIR>/node<id>_gt.json} & Ground truth del grafo.\\
\filepath{<WORK_DIR>/detection_cache/<image_stem>.json} & Detection e decoding memorizzati per immagine.\\
\filepath{<WORK_DIR>/multiview/node<id>/fused_graph.json} & Risultato principale.\\
\filepath{<WORK_DIR>/multiview/node<id>/scene_nodes.json} & Tracciabilità della ReID.\\
\filepath{<WORK_DIR>/multiview/node<id>/vggt/} & Artefatti geometrici per vista.\\
\filepath{<WORK_DIR>/multiview/node<id>/vggt_geometry.json} & Geometria aggregata, hint e cache VGGT.\\
\filepath{<WORK_DIR>/multiview/node<id>/fused_graph_plot.png} & Render prodotto da \filepath{visualize.py}.\\
\filepath{<WORK_DIR>/multiview/node<id>/standard_similarity.json} & Ranking del nodo rispetto agli standard.\\
\filepath{<WORK_DIR>/standard_comparison/} & Tabelle e report aggregati.\\
\filepath{<DEBUG_OUTPUT_DIR>/node<id>/} & Dump interni presenti soltanto con debug attivo.\\
\bottomrule
\end{longtable}

\section{Variabili di \filepath{.env}}

\begin{longtable}{P{0.34\textwidth}P{0.18\textwidth}P{0.35\textwidth}}
\toprule
\textbf{Variabile} & \textbf{Default} & \textbf{Significato}\\
\midrule\endhead
\envvar{IMAGES_DIR} & obbligatoria & Cartella immediata delle immagini.\\
\envvar{LABELS_EXT_DIR} & obbligatoria & Label estese per la valutazione.\\
\envvar{YOLO_MODEL_PATH} & obbligatoria & File dei pesi YOLO.\\
\envvar{ASSETS_DIR} & \filepath{./assets} & Mapping, porte e SVG Shurikode.\\
\envvar{ANNOTATIONS_DIR} & \filepath{./annotations} & Grafi ground truth.\\
\envvar{WORK_DIR} & \filepath{./output} & Root di risultati e cache.\\
\envvar{DEBUG_OUTPUT_ENABLED} & \code{false} & Abilita i dump interni.\\
\envvar{DEBUG_OUTPUT_DIR} & \filepath{./output/debug} & Root del debug.\\
\envvar{VGGT_EDGE_ENABLED} & \code{true} & Abilita geometria e hint VGGT.\\
\envvar{VGGT_GENERATE_IF_MISSING} & \code{true} & Genera i file \code{.npz} mancanti.\\
\envvar{VGGT_EMBEDDED_PYTHON} & \filepath{./python312/python.exe} & Runtime separato usato da VGGT.\\
\envvar{VGGT_MODEL_ID} & \code{facebook/VGGT-1B} & Identificativo del modello.\\
\envvar{VGGT_PREPROCESS_MODE} & \code{crop} & Pre-processing delle immagini VGGT.\\
\bottomrule
\end{longtable}

\section{Verifica minima del risultato}

Prima di considerare concluso un nodo:

\begin{itemize}
  \item verificare che \filepath{fused_graph.json} sia stato aggiornato dal run;
  \item controllare \code{is_connected}, il numero di nodi e il numero di archi;
  \item individuare archi con \code{forced=true} e valori
        \code{diameter_conflict=true};
  \item aprire \filepath{fused_graph_plot.png};
  \item conservare il log della valutazione;
  \item eseguire per ultimo l'eventuale export QGIS definitivo.
\end{itemize}
